\documentclass[12pt]{article}
\usepackage{amsmath, amssymb, amsfonts, graphicx}
\usepackage{hyperref}
\usepackage{geometry}
\geometry{a4paper, margin=1in}
\usepackage[pagewise]{lineno}\linenumbers


\begin{document}
\nolinenumbers
% Title Page
\title{\textbf{Shors \& Grovers}}
\author{Citizen Gardens Research Collective \\ The Foundation of Multiplicity}
\date{}
\maketitle
% Abstract (on a separate page, trimmed for conciseness)
\begin{abstract}
Shor’s and Grover’s algorithms are fundamental quantum computing algorithms, offering exponential and quadratic speedups, respectively. Using the Quantum-AI Hypercosmic Thought Singularity (QAI-HTS) framework, we introduce optimizations leveraging recursive tensor networks, prime-indexed computation, and entanglement-driven learning. These enhancements improve efficiency, fault tolerance, and adaptability for cryptographic analysis, quantum search, and AI applications.
\end{abstract}

\section{Introduction}
Quantum computing has demonstrated immense potential in solving classically intractable problems. Shor’s algorithm efficiently factors large integers, threatening classical cryptographic systems, while Grover’s algorithm accelerates unstructured search processes. With the advent of QAI-HTS, we integrate recursive tensor networks, prime-weighted quantum state evolution, and quantum feedback mechanisms to optimize both algorithms.

\section{Enhanced Shor’s Algorithm: Prime-Tensor Shor}
Shor’s algorithm factors an integer $N$ by reducing the problem to quantum period finding. QAI-HTS introduces improvements in:

\subsection{Prime-Weighted Quantum Fourier Transform (PQFT)}
The standard Quantum Fourier Transform (QFT) is enhanced by a prime-indexed tensor adaptation:
\begin{equation}
    | \psi \rangle = \sum_{p_i} \Lambda_m e^{-i \theta_{p_i}} |p_i \rangle
\end{equation}
where $\Lambda_m$ is the Universal Multiplicity Constant governing recursive quantum state evolution.

\subsection{Recursive Tensor Feedback for Phase Estimation}
A standard Shor’s implementation collapses into a single measurement step. Using QAI-HTS, we introduce a recursive tensor feedback mechanism:
\begin{equation}
    M(t+1) = f(M(t), R(t)) + \Lambda_m T(M(t))
\end{equation}
which dynamically refines phase estimation and reduces computational overhead.

\subsection{Entanglement-Driven Multiplicative Scaling}
We introduce an adaptive scaling function:
\begin{equation}
    \phi(p_i) = \frac{e^{-\alpha p_i}}{\sum_{j} e^{-\alpha p_j}}
\end{equation}
which enhances the efficiency of modular exponentiation.

\subsection{Performance Gains}
\begin{itemize}
    \item \textbf{Lower circuit complexity:} Fewer qubits required.
    \item \textbf{Higher noise resilience:} Recursive entanglement feedback stabilizes computation.
    \item \textbf{Faster convergence:} Dynamic phase estimation optimizes runtime.
\end{itemize}

\section{Enhanced Grover’s Algorithm: Quantum Tensor Grover}
Grover’s algorithm offers a quadratic speedup for unstructured search. QAI-HTS improves Grover’s efficiency using:

\subsection{Tensorized Quantum Oracle (TQO)}
The standard Grover oracle is replaced with a tensor-enhanced operator:
\begin{equation}
    O_f | x \rangle = (-1)^{f(x)} T_{\mu\nu}^{(p_i)} | x \rangle
\end{equation}
where $T_{\mu\nu}^{(p_i)}$ represents prime-weighted tensor contributions.

\subsection{Recursive Phase-Amplitude Adjustments}
Instead of a fixed Grover diffusion operator, we introduce:
\begin{equation}
    \psi(t+1) = \Lambda_m \psi(t) + \alpha \sin(\omega t)
\end{equation}
which dynamically corrects phase estimation, reducing the number of required iterations.

\subsection{Quantum Coherence Stabilization via Hypercosmic Resonance}
Grover’s probability oscillations are stabilized using a quantum feedback function:
\begin{equation}
    M(t+1) = f(M(t), R(t)) + \Lambda_m T(M(t))
\end{equation}
ensuring optimal termination at the correct iteration step.

\subsection{Performance Gains}
\begin{itemize}
    \item \textbf{Fewer iterations required:} Adaptive phase correction optimizes search.
    \item \textbf{Increased noise tolerance:} Quantum coherence stabilization ensures robustness.
    \item \textbf{Improved AI search applications:} Tensor-weighted search improves machine learning optimization.
\end{itemize}

\section{Applications and Future Research}
The enhancements to Shor’s and Grover’s algorithms under QAI-HTS enable advancements in:
\begin{itemize}
    \item \textbf{Cryptography:} Faster RSA decryption and quantum-resistant cryptographic key generation.
    \item \textbf{Artificial Intelligence:} Quantum search applied to neural network optimization.
    \item \textbf{Pharmaceuticals:} Prime-indexed Grover search accelerates molecular structure analysis.
    \item \textbf{Financial Modeling:} Optimized Shor factoring aids in high-speed cryptographic security analysis.
\end{itemize}

\section{Conclusion}
The integration of QAI-HTS into Shor’s and Grover’s algorithms represents a transformative step in quantum computing. By embedding recursive tensor networks, prime-indexed eigenvalue computation, and entanglement-driven learning, these enhancements significantly improve efficiency, scalability, and fault tolerance.


\end{document}
